\documentclass[12pt,a4paper]{report}

% ── Packages ──────────────────────────────────────────────────────────────────
\usepackage[utf8]{inputenc}
\usepackage[T1]{fontenc}
\usepackage[english]{babel}
\usepackage{geometry}
\usepackage{setspace}
\usepackage{titlesec}
\usepackage{hyperref}
\usepackage{booktabs}
\usepackage{array}
\usepackage{enumitem}
\usepackage{parskip}
\usepackage{graphicx}
\usepackage{fancyhdr}
\usepackage{caption}
\usepackage{microtype}
\usepackage{lmodern}
\usepackage{csquotes}
\usepackage{float}
\usepackage{xurl}

% ── Page Geometry ──────────────────────────────────────────────────────────────
\geometry{
  a4paper,
  top=2.5cm,
  bottom=2.5cm,
  left=3cm,
  right=2.5cm
}

% ── Line Spacing ───────────────────────────────────────────────────────────────
\onehalfspacing

% ── Header / Footer ───────────────────────────────────────────────────────────
\pagestyle{fancy}
\fancyhf{}
\rhead{\small AI-Driven XDR Simulation System}
\lhead{\small Diploma Thesis -- Preliminary Report}
\cfoot{\thepage}
\renewcommand{\headrulewidth}{0.4pt}

% ── Section Title Formatting ──────────────────────────────────────────────────
\titleformat{\chapter}[hang]{\Large\bfseries}{Chapter \thechapter.}{0.5em}{}
\titleformat{\section}[hang]{\large\bfseries}{\thesection}{0.5em}{}
\titleformat{\subsection}[hang]{\normalsize\bfseries}{\thesubsection}{0.5em}{}
\titlespacing*{\chapter}{0pt}{20pt}{15pt}
\titlespacing*{\section}{0pt}{14pt}{8pt}
\titlespacing*{\subsection}{0pt}{10pt}{6pt}

% ── Hyperref Setup ────────────────────────────────────────────────────────────
\hypersetup{
  colorlinks=true,
  linkcolor=blue,
  urlcolor=blue,
  citecolor=blue,
  pdftitle={AI-Driven XDR Simulation System},
  pdfauthor={Darling Selita}
}

% ── Document ──────────────────────────────────────────────────────────────────
\begin{document}

% ── Title Page ────────────────────────────────────────────────────────────────
\begin{titlepage}
  \centering
  \vspace*{2cm}
  {\LARGE\bfseries Diploma Thesis -- Preliminary Report\par}
  \vspace{1.5cm}
  {\Large\bfseries AI-Driven XDR (Extended Detection and Response) Simulation System\par}
  \vspace{2cm}
  {\large Darling Selita\par}
  \vspace{1cm}
  {\normalsize 2026\par}
  \vfill
\end{titlepage}

% ── Abstract (Albanian) ───────────────────────────────────────────────────────
\chapter*{Abstrakt}
\addcontentsline{toc}{chapter}{Abstrakt}

Teknologjia ndodhet në një nga periudhat më kulminante të saj, duke sjellë zhvillim dhe
përmirësim të vazhdueshëm në të gjitha fushat e jetës sonë të përditshme. Por, me rritjen e
varësisë ndaj sistemeve digjitale apo të automatizimeve, duhet të merret parasysh dhe siguria e
tyre, në mënyrë që këto sisteme të ofrojnë shërbim të mirë dhe të vazhdueshëm, duke e bërë
kështu, sigurinë e të dhënave një sfidë të madhe për mbarëvajtjen.

Sistemet tradicionale të përdorura, si SIEM, EDR, fokusohen kryesisht tek parandalimi i këtyre
sulmeve, duke skanuar herë pas herë log-et e sistemit, në mënyrë që të kuptojë se kur një
aktivitet është normal apo i dyshimtë, duke kërkuar gjithashtu një ndërhyrje njerëzore për të
kuptuar se çfarë ka ndodhur dhe si duhet përgjigjur. Këtu hyn në punë XDR, një sistem i cili
synon të integrojë fjalën e fundit të teknologjisë, Inteligjencën Artificiale, dhe nëpërmjet një
baze të qendërzuar të të dhënave nga ku të gjithë kompanitë, të cilat shesin shërbime për
mbrojtjen kibernetike, jo vetëm marrin të dhëna por edhe e përditësojnë këtë bazë të dhënash
në rast të një virusi apo anomalie të re të shfaqur në sistemet e tyre përkatëse, të automatizojë
të gjithë procesin e kundërpërgjigjes ndaj sulmeve apo incidenteve.

Kjo temë propozon zhvillimin e një sistemi simulues XDR, të bazuar në Inteligjencë Artificiale,
i cili do ketë si qëllim jo vetëm analizimin e log-eve të sistemit, por edhe marrjen e një
vendimi autonom për kundërpërgjigje. Ky sistem XDR do të përdorë Machine Learning, për të
trajnuar si të identifikohen sjellje jo-normale, si dhe të simulohen veprime për të izoluar dhe
bllokuar një proces apo një instancë të dyshimtë.

Qëllimi përfundimtar, përtej simulimit, do të jetë vlerësimi i efikasitetit dhe funksionalitetit të
kësaj mënyre të të punuarit, për të vlerësuar më pas gjithashtu dhe besueshmërinë, saktësinë
dhe sigurinë e vendimmarrjes së këtij sistemi apo çdo lloj sistemi inteligjent.

\medskip
\noindent\textbf{Fjalë kyçe:} \textit{Inteligjencë Artificiale, Siguri e të dhënave, SIEM, EDR,
XDR, Machine Learning, Analizimi i log-eve, Sisteme digjitale, Vendimmarrje autonome.}

% ── Abstract (English) ───────────────────────────────────────────────────────
\chapter*{Abstract}
\addcontentsline{toc}{chapter}{Abstract}

Technology is in one of its most culminating periods, bringing continuous development and
improvement in all areas of our daily lives. But, with the increasing dependence on digital
systems or automation, their security must also be taken into account, so that these systems
provide good and continuous service, thus making data security a major challenge for the
well-being.

Traditional systems used, such as SIEM and EDR, focus mainly on preventing attacks by scanning
system logs periodically, in order to understand when an activity is normal or suspicious, also
requiring human intervention to understand what happened and how to respond. This is where XDR
comes into play---a system that aims to integrate the latest technology, Artificial Intelligence,
and through a centralized database from which all companies that sell cyber protection services
not only receive data but also update this database in the event of a new virus or anomaly
appearing in their respective systems, to automate the entire process of responding to attacks or
incidents.

This thesis proposes the development of an XDR simulation system, based on Artificial
Intelligence, which will aim not only to analyse system logs, but also to make an autonomous
decision to respond. This XDR system will use Machine Learning to train how to identify abnormal
behaviour, as well as simulate actions to isolate and block a suspicious process or instance.

The ultimate goal, beyond simulation, will be to evaluate the efficiency and functionality of
this way of working, to then also assess the reliability, accuracy, and security of
decision-making of this system or any type of intelligent system.

\medskip
\noindent\textbf{Keywords:} \textit{Artificial Intelligence, Data Security, SIEM, EDR, XDR,
Machine Learning, Log Analysis, Digital Systems, Autonomous Decision-Making.}

% ── Table of Contents ─────────────────────────────────────────────────────────
\tableofcontents
\clearpage

\listoffigures
\clearpage

% ══════════════════════════════════════════════════════════════════════════════
\chapter{Introduction}
% ══════════════════════════════════════════════════════════════════════════════

In modern times, technology hasn't been stopping one bit, and has kept climbing the steps of
evolution rapidly, therefore transforming our everyday life, and keeping us in touch with it for
our day-to-day activities. Not only automated systems, but also digital systems and big complex
infrastructures, that tech companies use to operate their work environment on, have been deeply
embedded with data collection \& processing, communication and service delivery. However, these
advancements come with their own problems, one of which---if not the most important one---is the
cybersecurity threats that these infrastructures face every day.

As more organizations, not only in IT or Cybersecurity especially, depend on these customs,
there is also a big need of insurance for these services, not only that they function well, but
also that their service is uninterrupted, and the data being used is correct and not always
tampered with, as well as available to the instances who need it. Therefore, Cybersecurity has
become a crucial component to keeping these systems and infrastructures running at all times and
without errors.

% ──────────────────────────────────────────────────────────────────────────────
\section{Problem Statement}
% ──────────────────────────────────────────────────────────────────────────────

Cybersecurity solutions that are mostly used as of now are SIEMs (Security Information Event
Management) and EDRs (Endpoint Detection Services), to monitor system logs and detect potential
suspicious activity, or potential threats. This is done continuously, thus giving these systems a
hand when comparing normal and abnormal activity.

However, despite their good traits and effectiveness, these solutions present limitations for the
current rate of threats that systems face. They sometimes require cyber threat analysts to be
involved in the decision-making process, in order to better analyse alerts and investigate these
incidents, thus leading to a considerable delay in response time and an increase in the amount of
work from these human experts. But as mentioned above, these limitations are currently present,
and if we're thinking ahead, there needs to be consideration for the advancements of these
threats, of the ways they infiltrate the system, and how hard they hit.

% ──────────────────────────────────────────────────────────────────────────────
\section{Motivation}
% ──────────────────────────────────────────────────────────────────────────────

With the problem being apparent, and having a considerable impact on security, XDR is not only a
solution to the problem, but also a stepping stone for the next era of incident response systems,
coupled with the advancements of AI, offering a better range of support and detection abilities,
as well as more automated responses that greatly reduce alert lateness. It is also important to
mention that the data that these systems will be fed with will not be only new data, or only old
data, but a variation that will help the AI model to have a better sense of direction when
determining whether a log is suspicious or not. The costs of developing these solutions also
greatly overextend the cost of human labour, making it a better choice for the future.

% ──────────────────────────────────────────────────────────────────────────────
\section{Research Questions}
% ──────────────────────────────────────────────────────────────────────────────

The main goal of this thesis is to find out how effective AI and powerful ML algorithms are when
integrated in real-world Extended Detection and Response Systems (XDRs). The concrete questions
raised in this thesis are:

\begin{itemize}
  \item[\textbf{RQ1:}] Can an Artificial Intelligence-driven XDR System accurately identify
  suspicious logs or events using a centralized telemetry system along with continuous log
  monitoring and analysis?

  \item[\textbf{RQ2:}] How much can Machine Learning algorithms improve threat detection, event
  correlation and risk mitigation?

  \item[\textbf{RQ3:}] Is it possible to simulate an almost fully automated incident response
  mechanism with an AI + ML Engine, while maintaining the level of reliability and accuracy that
  normal XDR solutions offer?

  \item[\textbf{RQ4:}] How effective is intrusion detection and event correlation in early
  mitigation of sequential or multi-scale attacks?
\end{itemize}

% ──────────────────────────────────────────────────────────────────────────────
\section{Aim and Objectives}
% ──────────────────────────────────────────────────────────────────────────────

The main objective of this thesis is to design and develop a miniature version of an XDR system,
based on Artificial Intelligence and Machine Learning. To destructurize this objective, the
following course of action is:

\begin{itemize}
  \item Analyse existing Cybersecurity solutions (SIEM, EDR, SOAR).
  \item Go through Machine Learning techniques to find what is best suited to detecting
  suspicious activity.
  \item Design a system capable of pulling and analysing data from a data source, as well as
  develop an AI model to detect anomalies and return the appropriate response such as target
  isolation, blocking and process killing.
  \item Evaluate the metrics of this system, to determine if it is reliable and safe to be
  applied on a much larger scale.
\end{itemize}

% ──────────────────────────────────────────────────────────────────────────────
\section{Methodological Review}
% ──────────────────────────────────────────────────────────────────────────────

The first part of the study involves an overview of existing Cybersecurity Solutions, AI-based
detection techniques and XDR systems, found in academic reviews and research as well as industry
standards such as NIST or ISO.

The next part will cover the design and development of the proposed XDR system, with the focus
of improving the existing ones. Based on the review, there will need to be an assessment of the
correct technologies and tools needed to develop this system, to not only be functional, but if
successful, even scalable, with focus on data processing quality and model selection.

This step will also involve finding the appropriate datasets to train the model with, once
selected, to then apply ML algorithms to detect suspicious behaviour and simulate needed
responses.

After all of that, the last step will be to thoroughly analyse the system, using performance
metrics such as detection accuracy, response effectiveness towards solving the problem,
reliability and availability. This will also help when drawing conclusions, in order to determine
any setbacks that will need further research or development.

% ══════════════════════════════════════════════════════════════════════════════
\chapter{Literature Review}
% ══════════════════════════════════════════════════════════════════════════════

This chapter represents a view of the existing literature for Cybersecurity solutions and defence
systems, in regard to the evolution of complexity and intensity of cyber threats. As the systems
that people use every day continue to scale based on their needs, the requirement for good
security to prevent unauthorized access to these systems becomes more crucial.

The primary parts will include an analysis of SIEM, EDR and SOAR systems, how they work and how
effective they are with the current range of cyber threats, followed by the turning points that
an XDR system will offer to solve existing problems or fatigue in regard to anomaly detection and
resolution.

This will be done by also expanding on the role of Artificial Intelligence and Machine Learning
in enhancing such systems and building upon them, to improve these mechanics.

Finally, the chapter will identify limitations and problems of current approaches, that also show
the research gap, which in turn requires the innovation and development of systems such as XDRs.

% ──────────────────────────────────────────────────────────────────────────────
\section{Cybersecurity Basics and Threat Environment}
% ──────────────────────────────────────────────────────────────────────────────

Big companies and organizations rely heavily on digital systems or platforms, which they use to
offer their users their services. In the background, all these systems and platforms are run
through data collection and interpretation subsystems, applied for communication, data storage
and service delivery (CDNs). However, this also brings upon the issue of security, integrity and
availability of these subsystems, as they are crucial to maintaining these platforms operational
and running.

Cyber threats have evolved almost at the same pace as technology advancements, thus making it
difficult to maintain a level of security within the infrastructure. According to the MITRE
ATT\&CK framework~\cite{ref1}, phishing (T1566) remains one of the most predominant attacks,
occurring daily and directed to every endpoint of the infrastructure. The low cost -- high reward
matrix makes it a very favourable decision for the attacker to use and establish a persistent
connection with the victim, unknowingly.

However, behind these statistics lies a vast amount of information regarding techniques, tactics
and procedures, that helps cyber threat analysts better understand these attacks and how to
interpret the alerts that a SIEM or an EDR show.

The NIST framework~\cite{ref2} gives a clear overview of what an event is, how events are
categorized, and how to know what to be aware of. It also emphasizes the need for incident
response mechanisms, including protecting data (personal or business), as well as keeping
integrity of the systems. The proposed XDR has the scope of doing so in a shorter amount of
time, also keeping precision and correctness during the decision-making process.

\begin{figure}[htbp]
  \centering
   \includegraphics[width=0.8\textwidth]{fig01_nist_lifecycle.png}
  \caption{Incident Response Lifecycle (NIST.SP.800-61r2)}
  \label{fig:nist_lifecycle}
\end{figure}

The other problem that needs to be addressed is the increase of better set-out attacks, that
require a more complex course of action needed to be taken in order to defend and respond.
Even with the use of the MITRE ATT\&CK framework, which documents new methods, the responses
are still delayed and slow due to interpretation or understanding of issues.

% ──────────────────────────────────────────────────────────────────────────────
\section{Traditional Detection Systems: SIEM and EDR --- Procedure \& Limitations}
% ──────────────────────────────────────────────────────────────────────────────

Corporate security monitoring and compliance projects have traditionally started with Security
Information and Event Management (SIEM) solutions. The main purpose of these systems is to
collect log data and perform continuous checks, as well as work in parallel with other
cybersecurity solutions like IDS (Intrusion Detection Systems) and IPS (Intrusion Prevention
Systems) from which log data is also taken. Most SIEMs use a centralized language, translating
the logs from their native language into one unified language~\cite{ref3}.

\begin{figure}[H]
  \centering
   \includegraphics[width=0.8\textwidth]{fig2_threat_detection_evolution.png}
  \caption{Evolution of Threat Detection~\cite{ref3}}
  \label{fig:evolution_detection}
\end{figure}

SIEMs have been used for years, but as noted before, the technology advancements have made it
even more difficult to keep up with new and sophisticated threats. SIEMs remain a good choice for
overall management, coupled with a very insightful dashboard in most cases, but when a fast-paced
attack occurs, these systems may not be able to appropriately deliver a response in the required
amount of time, resulting in loss of data, resources and time.

On the other hand, Endpoint Detection \& Response (EDR) systems are concentrated on identifying
and reacting to the unique threats discovered on endpoints (end-users or workers of a
company/organization), giving a much more focused insight into endpoint activities, monitoring,
detection and appropriate response. Due to better visibility from an EDR, cyber threat analysts
are more prone to using it on a day-to-day basis~\cite{ref4}.

However, EDR offers a more centralized view of the data but is limited to a single endpoint, or
a small number of endpoints, making it harder to manage a large number of endpoints in real time.
This is the case for most organizations, which have many devices---not just computers, but
servers and CDNs as well---thus bringing out the limitations of this technology.

\begin{figure}[htbp]
  \centering
   \includegraphics[width=0.8\textwidth]{fig3_edr_methodology.png}
  \caption{EDR Work Methodology~\cite{ref4}}
  \label{fig:edr_methodology}
\end{figure}

% ──────────────────────────────────────────────────────────────────────────────
\section{Extended Detection and Response (XDR)}
% ──────────────────────────────────────────────────────────────────────────────

To reiterate the previously mentioned problem: the cyber threats of today are much more complex
and scaled, and with traditional systems you can only achieve so much protection against them. On
these grounds, with the recent developments of AI and cross-layer visibility, a new security
approach has been established---the Extended Detection and Response system, XDR---showing a
significant improvement in threat visibility, detection accuracy and emergence when responding to
an incident.

An XDR is not the kind of system that should be deployed everywhere, as an EDR and SIEM can work
together just fine if the architecture is not complex and the system is not easily accessible.
An XDR handles information from various sources, much like a SIEM, collects all these logs, and
then initiates a specific response based on what the logs convey. Sources of information include
endpoints, networks, cloud architectures and so on. The effectiveness of this system is higher
than an EDR or a SIEM by themselves, because of the correlation between different sources of
information and layers, making it have a better understanding of security events and enabling
faster detection of threats.

According to Microsoft, an XDR solution is designed to unify the work of two or more separate
systems into one, reducing the complexity of having to manage and maintain different security
tools~\cite{ref5}.

One of the advantages of these systems is the ability to automate their actions and responses
based on analytics and previous data, which they take through large and public data sets, and
then use to continuously update and train their machine learning algorithms~\cite{ref6}. This
way, things like isolating compromised endpoints, blocking malicious processes or events,
restricting access to the network, or even fully blocking a specific endpoint are things that now
do not require any manual intervention or predefined rules.

Another important source of information for XDR systems is the MITRE ATT\&CK framework, a
globally accessible knowledge base of adversary tactics and techniques, based on real-world
observations~\cite{ref7}. It provides structured information about TTPs (Techniques, Tactics and
Procedures) that the XDR can use to classify and correlate malicious behaviour throughout system
nodes and layers of the architecture.

Despite this, it should be taken into account that an XDR solution does not automatically fix any
shortcomings that an organization might have. It is still an imperfect concept, with most XDR
systems still relying on predefined sets of rules, having issues with the correlation between
different logs from different sources due to training limitations, and most importantly, still
potentially requiring human intervention in the case of unique alerts or final response decisions.

As cybersecurity threats become more sophisticated, the second generation of XDR systems began
to incorporate AI and machine learning technologies altogether. These advancements allowed for
better anomaly detection, predictive analytics and automated response capabilities. XDR systems
could now learn from data, and directly apply the knowledge learned in response to an
incident~\cite{ref6}.

% ──────────────────────────────────────────────────────────────────────────────
\section{Artificial Intelligence and Machine Learning in Cybersecurity}
% ──────────────────────────────────────────────────────────────────────────────

The increasing difficulty of having to face complex and frequent cyberattacks has raised serious
doubts about the rule-based, traditional security systems. As a result, with the second
generation of XDRs, AI and Machine Learning technologies have been embraced as much as possible
in an effort to raise the system's awareness to patterns and malicious intents, through analysing
large volumes of past data (datasets) and present data (logs arriving to the system through
endpoints, layers, network, etc).

The integration of Artificial Intelligence and Machine Learning into cybersecurity has driven a
transformational shift, significantly enhancing the ability to detect, respond to and mitigate
complex cyber threats~\cite{ref8}, with AI becoming a cornerstone in advancing cyber security,
offering innovative solutions to detect and mitigate cyber threats~\cite{ref9}.

\begin{figure}[htbp]
  \centering
   \includegraphics[width=0.8\textwidth]{fig4_ai_ml_techniques.png}
  \caption{The key AI and ML techniques applied in cybersecurity~\cite{ref19_fig}}
  \label{fig:ai_ml_techniques}
\end{figure}

AI and ML are recognized as effective tools that enhance cybersecurity systems and allow for
real-time analytics and decision-making while supporting applications such as intrusion detection,
malware detection and network security~\cite{ref10}. AI gives systems the ability to think and
reason just like humans, incorporating pattern-recognition and problem-solving. In the vast world
of machine learning techniques used in cybersecurity, three key approaches stand out: supervised,
unsupervised and reinforcement learning~\cite{ref11}.

\begin{figure}[htbp]
  \centering
   \includegraphics[width=0.7\textwidth]{fig5_zero_day_detection.png}
  \caption{Zero-Day Detection Rate by Approach~\cite{ref11}}
  \label{fig:zero_day_detection}
\end{figure}

% ──────────────────────────────────────────────
\subsection{Learning Methods Overview}
% ──────────────────────────────────────────────

In concrete terms, Supervised Learning is designed to work with the use of labels, giving the
XDR system a clear path of how to detect and correlate between irregular or suspicious patterns.
Unsupervised Learning is used without specific labels, giving the system a different
``perspective'' of anomalies or unusual patterns. Lastly, reinforcement learning is an approach
where agents essentially learn how to spot and contain malware by engaging in trial and error
within simulated environments. The learning agent decides whether to classify, quarantine or
allow a sample, adjusting its strategy based on the outcomes~\cite{ref12}.

Supervised learning shines in controlled settings, unsupervised learning in spotting zero-day
attacks, while reinforcement learning shines in adapting to ever-changing, adversarial
situations~\cite{ref12}.

This gives the system the possibility to handle any incoming situation and make sure to properly
mitigate an incoming issue. One of the most used features of security systems is Anomaly
Detection, focusing on identifying deviations from normal patterns (normal system behaviour),
allowing organizations to detect potentially malicious activities before dealing any significant
damage to the core structure. These detection sub-systems are trained through the use of large
datasets, one of which is the CICIDS2017 dataset, developed by the Canadian Institute for
Cybersecurity~\cite{ref13}.

To reiterate, these advancements in AI and ML are good, but also present some issues and
challenges when applying them in real-world solutions. First, the issue of diversity between logs
and information that systems are trained with. If the data is not diverse enough, the system will
not be able to see through unique or zero-day attacks. Conversely, if the dataset is more diverse
than usual, the recognition of daily suspicious patterns becomes hard to see and prolongs
automated responses for small anomalies.

It is also worth noting that AI allows new vulnerabilities in cybersecurity. Criminals,
adversaries, and malicious actors use AI and other technologies to circumvent detection systems
and automate attacks with minimal human intervention. The dual-use nature of AI raises major
risks in the areas of cybersecurity, privacy, and public trust~\cite{ref14}.

\begin{figure}[htbp]
  \centering
   \includegraphics[width=0.85\textwidth]{fig6_ai_threats_taxonomy.png}
  \caption{Taxonomy of AI-Driven Cybersecurity Threats and Corresponding Defensive
  Strategies~\cite{ref14}}
  \label{fig:ai_threats_taxonomy}
\end{figure}

% ──────────────────────────────────────────────────────────────────────────────
\section{Log Analysis and Automated Response Systems}
% ──────────────────────────────────────────────────────────────────────────────

The analysis of logs for anomalies is an important research topic with practical importance in
the field of failure analysis and security threat detection. Logs are generated by software-driven
applications running on systems or devices to provide information to aid software developers and
system engineers with their analysis of system state and condition~\cite{ref15}. Logs are
required to be generated by any software in order for the security systems to monitor ongoing
processes and maintain a level of security, by detecting any unusual patterns or investigating
incidents based on suspicious activity. Due to the increasing volume and complexity of this
generated data, effective log management and analysis is being prioritized for development, as it
is one of the most effective ways to improve threat detection.

% ──────────────────────────────────────────────
\subsection{Currently Well-Established Security Systems \& Innovations}
% ──────────────────────────────────────────────

Currently, the most utilized system for gathering logs and processing extracted information are
SIEMs, with notable mentions such as Microsoft Sentinel and Splunk. These systems collect,
aggregate and correlate log data, empowering network defenders to monitor activity and uncover
advanced cyberthreats~\cite{ref16}. The correlation of events from multiple platforms helps
detect sequences of related events and highlight abnormalities that may not be immediately visible
to a human when manually checking the logs~\cite{ref17}.

Setting SIEMs aside, most organizations that require cybersecurity services rely heavily on
automatic response systems, such as SOAR (Security Orchestration, Automation, Response). These
systems not only automate their responses based on the detected pattern of abnormalities, but
also co-operate with other security tools to implement a specific course of action, coordinated
throughout multiple tools, against any suspicious activity that is detected~\cite{ref16}. The
responses can vary from normal ones, such as alert generation or blocking suspicious IP addresses,
to complex ones, such as isolating a compromised device or terminating a malicious process.

However, these security systems still work based on predefined static rules, which limit their
ability to evolve and handle complex unknown threats, with humans having to check every alert
and decide whether the appropriate response is being used by the system.

\begin{figure}[htbp]
  \centering
   \includegraphics[width=0.75\textwidth]{fig7_ai_soar.png}
  \caption{Integration of AI/ML into SOAR platforms}
  \label{fig:ai_soar}
\end{figure}

These limitations, along with the evolving cycle that threats have gone through---from virus
propagation and phishing, to large-scale attacks that blend stealth, automation and
persistence~\cite{ref18}---have shifted the cybersecurity field towards the AI approach, where
more intelligent systems can be developed to effectively detect, respond to, and mitigate even
complex cyber-attacks.

% ──────────────────────────────────────────────────────────────────────────────
\section{Autonomous AI and Research Gap}
% ──────────────────────────────────────────────────────────────────────────────

To reinforce what has been established in the previous sections, the evolution of cyber threats
has increased the demand for protection against such threats, but with a difference from the
traditional approach: more automation and intelligent adaptation to cyber threats (both known
ones and zero-day) and minimal human intervention. While traditional security systems are
momentarily still a secure choice for any organization or state infrastructure, with these
threats advancing and with the way that organizations are scaling towards much bigger
infrastructures, human intervention will become much harder and inefficient when faced with a
multi-scale attack.

% ──────────────────────────────────────────────
\subsection{Autonomous AI Systems \& Agentic AI Security Solutions --- an Overview}
% ──────────────────────────────────────────────

Autonomous AI systems are designed for the simulation of decision-making processes, by firstly
analysing and processing all the data given to them, determining visible patterns, and then
deciding on the risks of different actions that they initiate. In security solutions, these
sub-systems can work as a combination of different security tools, in order to detect and
determine whether an attack is coming, and then execute response actions with minimal to no
human intervention. One of the key advantages of these AI/ML systems is their ability to adapt
to incoming threats, then develop specific patterns of response actions to deal with them,
therefore evolving alongside the threats they are fighting against~\cite{ref19b}.

\begin{figure}[htbp]
  \centering
   \includegraphics[width=0.75\textwidth]{fig7b_adaptive_defence.png}
  \caption{Adaptive defence system for cybersecurity~\cite{ref20}}
  \label{fig:adaptive_defence}
\end{figure}

The picture above gives an illustration of how an autonomous AI system currently works, divided
into components for a better visual representation. First, we have the data capturing tools,
which are usually watching and monitoring all the available events in the network for logs, such
as Winpcap (Windows) and Libpcap (Linux)~\cite{ref20}. After that, data preprocessing tools
consider all these events, and try to separate the ones that they have already learned from or
mastered, filtering them out, and leaving only the ones that are unknown. The next process, of
feature extraction, as its name suggests, extracts features that make this event unique and
suspicious, which then pass it to analysis engines, in order to form a plan of response. After
responding, the next layer of security activates, the Intrusion Detection System (IDS)~\cite{ref20}.

\begin{figure}[htbp]
  \centering
   \includegraphics[width=0.7\textwidth]{fig8_agentic_ai.png}
  \caption{Agentic AI anatomy~\cite{ref23}}
  \label{fig:agentic_ai}
\end{figure}

Another example of an innovative technology that integrates AI and LLMs in cybersecurity is
Agentic AI. Nowadays, most large language models serve a key role in data preprocessing, and
response generation in regard to the question or problem posed. They are able to process very
large amounts of data, recognize patterns, learn from them and identify any patterns that can
result in malicious activity~\cite{ref22}. Moreover, its capability of reading rules and
guidelines, along with detailed analysis about certain attacks, helps analysts create a more
structured plan of response, even without executing that response itself~\cite{ref22}.

% ──────────────────────────────────────────────
\subsection{Current Problems of AI \& ML Security Solution Approach}
% ──────────────────────────────────────────────

\begin{figure}[htbp]
  \centering
   \includegraphics[width=0.75\textwidth]{fig9_ai_ml_security_automation.png}
  \caption{AI/ML in Security Automation~\cite{ref21}}
  \label{fig:ai_ml_security_automation}
\end{figure}

The aim of this thesis is not only to support the use of Artificial Intelligence in cybersecurity,
but also to show the gaps that are currently present, and how the integration of Cybersecurity,
AI and Machine Learning will help fill them. One challenge that we're faced with in this approach
is the low amount of XDR-focused simulation systems. Existing studies focus on sole features,
such as Machine Learning patterns for detecting anomalies, or rule-based automation using AI,
therefore not creating a single instance application which integrates both these functions and
gives more functionality in one tool~\cite{ref21}.

Rule-based systems can also be integrated with other security systems, such as intrusion
detection systems and SIEM systems~\cite{ref21}, thus making a more complete security solution
that covers all the required functionalities.

Another issue is the problem of trust in the accuracy of threat detection and response decisions.
This is a persisting issue, present even in current deployed security solutions, where these
systems lack the ability to adapt to dynamically evolving threats~\cite{ref22}, and therefore
cannot be fully autonomous.

% ──────────────────────────────────────────────────────────────────────────────
\section{Ethical and Operational Challenges of AI-Driven Cybersecurity Systems}
% ──────────────────────────────────────────────────────────────────────────────

Now that we have established the importance of integrating AI in cybersecurity, with two of the
most important reasons being the speed with which threats are evolving day by day and the
improvements that AI has shown for the effectiveness and efficiency of threat detection and
response~\cite{ref24}, it is also worth noting that this integration should occur by taking into
account several ethical and operational issues. The need for more intelligent systems should also
bring the need for concerns over the reliability and transparency that these proposed security
solutions offer, in comparison with conventional systems.

\begin{figure}[htbp]
  \centering
   
\includegraphics[width=0.5\textwidth]{fig10_chatgpt_disclaimer.png}
  \caption{ChatGPT User Interface Disclaimer}
  \label{fig:chatgpt_disclaimer}
\end{figure}

% ──────────────────────────────────────────────
\subsection{Reliability and Accuracy of AI Decisions}
% ──────────────────────────────────────────────

One of the main concerns surrounding the current state of AI is the reliability that it offers
to its users. From simple things like searching for information, asking a question or posing a
problem, to complex situations such as detecting suspicious patterns and anomalies, there is
always a label that any AI software uses to inform the entity using it that AI results may not
be correct, and need further verification. Machine Learning models depend heavily on training
data, pre-processed by LLMs, and issues such as data quality, bias, and incomplete or poisoned
datasets (an attack technique used in Cybersecurity) can negatively affect detection accuracy and
system performance~\cite{ref21}.

For a normal user, this can only be a minor inconvenience, as the sources of information are
infinite nowadays. However, for a cybersecurity system facing a cyber-attack, one wrong mistake
can compromise entire systems or even the whole infrastructure of the organization, causing
full-service disruption. Therefore, creating a set of standards and best practices to ensure
high detection accuracy, and to minimize any decision errors, is a main priority.

% ──────────────────────────────────────────────
\subsection{Transparency and Explainability}
% ──────────────────────────────────────────────

Another critical issue that needs to be addressed is the transparency of a model's reasoning for
its actions and responses. Most AI large language models work in the form of Black Box systems,
which means that the process they go through to evaluate the best course of action is not easily
interpretable~\cite{ref25}, making it harder for humans to understand how the model arrived at a
certain point, and further validate its response. An analyst must make sure to understand why an
alert was triggered, how the pattern was discovered, and if the action set as a response is
correct or not. This is also a very important feature of the envisioned fully autonomous security
systems, as the decisions of these systems should be fully explainable to ensure trust and
accountability~\cite{ref25}.

\begin{figure}[htbp]
  \centering
   \includegraphics[width=0.75\textwidth]{fig11_xai_vs_blackbox.png}
  \caption{Explainable AI vs.\ Blackbox AI~\cite{ref26}}
  \label{fig:xai_blackbox}
\end{figure}

% ──────────────────────────────────────────────
\subsection{Security and Privacy Concerns (CIA Triad)}
% ──────────────────────────────────────────────

AI cybersecurity systems rely heavily on LLMs that process heavy amounts of data, be it data from
external sources (e.g.\ CICIDS 2017/2018) or internal sources such as the logs generated by the
events of the organisation infrastructure. Part of these internal logs might contain sensitive
information that needs to be handled carefully. The CIA triad (Confidentiality, Integrity and
Availability) forms the cornerstone of information security practice and theory, pointing out the
most important objectives that should be followed by these systems to guarantee a safe environment.

Furthermore, AI models as mentioned before might also be attacked through the datasets that they
use to learn from, using data poisoning~\cite{ref27}, that involves either of the following
actions:

\begin{itemize}
  \item Overfitting the data with the same type of information, thus making the model not
  adaptable to new threats.
  \item Incorrect logs that help create fake patterns that will be useless for detection or
  response action decision.
\end{itemize}

Another present issue is the source of information from where these events and logs are
generated. Second-generation XDR platforms quite often rely on different sources of information,
making it important to evaluate the validity of the information being used, so as to not increase
exposure to threats.

\begin{figure}[H]
  \centering
   \includegraphics[width=0.5\textwidth]{fig12_cia_triad.png}
  \caption{CIA Triad~\cite{ref28}}
  \label{fig:cia_triad}
\end{figure}

% ──────────────────────────────────────────────
\subsection{Importance for Next-Generation Security Systems}
% ──────────────────────────────────────────────

Even though presented with the current challenges, the field of AI, especially Agentic AI, is
evolving rapidly, making it possible to think that all these issues can be resolved. The envisioned
intelligent systems will provide a much better experience than the current ones, with more
autonomy in detection and decision-making. AI-driven security systems will cover multi-platform
security functionalities (detection, adaptation and automated response actions), forming the
foundation for fully automated and proactive system architectures~\cite{ref29}.

\begin{figure}[H]
  \centering
  \includegraphics[width=0.75\textwidth]{fig13_incident_response_lifecycle.png}
  \caption{Incident Response Lifecycle~\cite{ref29}}
  \label{fig:incident_response_lifecycle}
\end{figure}

% ──────────────────────────────────────────────────────────────────────────────
\section{XDR Systems and Compliance with Cybersecurity Frameworks}
% ──────────────────────────────────────────────────────────────────────────────

XDR + integrated AI systems are considered to be the next destination that cybersecurity is
heading to, considering how threats are evolving and organizations are continuously scaling their
infrastructure. XDR solutions really do offer compelling benefits; however, companies that are
already implementing or will implement them face significant challenges~\cite{ref30}. These
issues can be mitigated by implementing these solutions in compliance with specific standardized
cybersecurity frameworks.

The integration between these systems and cybersecurity frameworks helps identify threats earlier
by using well-known techniques, as well as creating a standardized set of procedures that should
be followed when responding to an incident. This section will further explain how XDR solutions
can work together with some of the most relevant frameworks: MITRE ATT\&CK, NIST and ISO/IEC
27001.

% ──────────────────────────────────────────────
\subsection{MITRE ATT\&CK Integration in XDR Systems}
% ──────────────────────────────────────────────

The MITRE ATT\&CK is one of the most widely used open-source frameworks in cybersecurity
solutions. It is a globally accessible knowledge base, containing documentation about real-world
threats, focusing on three key aspects of attacks: tactics, techniques and procedures
(TTPs)~\cite{ref31}.

This framework maps detected suspicious activities, reported in publicly accessed databases, to
known attacker behaviour, helping organizations detect and neutralize threats much faster. Not
only that, but XDR systems can ingest information from MITRE ATT\&CK, and use it to triage
alerts, enrich cyber threat intelligence from different sources and ultimately test the
effectiveness and accuracy of the incident response actions that it implements
autonomously~\cite{ref32}.

This framework also helps XDR systems to better understand cross-platform threats, such as
multi-stage intrusions that usually happen in broader networks, where the XDR is capable of
detecting suspicious patterns throughout all connected platforms.

It is also worth noting that the TTPs, characteristic of MITRE ATT\&CK and the way their
information is organized, make the use of this data for training AI models much more feasible.

\begin{figure}[H]
  \centering
   \includegraphics[width=0.8\textwidth]{fig14_mitre_attack.png}
  \caption{The MITRE ATT\&CK Framework~\cite{ref33}}
  \label{fig:mitre_attack}
\end{figure}

% ──────────────────────────────────────────────
\subsection{NIST Cybersecurity Framework and XDR}
% ──────────────────────────────────────────────

The NIST Cybersecurity Framework (CSF) is a unified framework that provides a structured approach
for managing a threat in a high-risk situation~\cite{ref34} using five core steps:

\begin{enumerate}
  \item Identify
  \item Protect
  \item Detect
  \item Respond
  \item Recover
\end{enumerate}

\begin{figure}[htbp]
  \centering
   \includegraphics[width=0.6\textwidth]{fig15_nist_framework.png}
  \caption{NIST Framework~\cite{ref34}}
  \label{fig:nist_framework}
\end{figure}

These steps are a functional dependency for XDR and AI security systems, which in turn provide
much better detection and visibility towards incoming threats, enabling a fast incident response
time. Using events from multiple points of the network and unifying data from these endpoints,
XDR systems enable real-time threat detection and automated response actions.

Furthermore, the NIST Cybersecurity Framework unifies the language for risk management, enabling
organizations to align security operations with business and compliance requirements~\cite{ref34}.

% ──────────────────────────────────────────────
\subsection{ISO/IEC 27001 and XDR Systems}
% ──────────────────────────────────────────────

ISO/IEC 27001 is a cybersecurity standard recognized internationally for Information Security
Management Systems (ISMS). It provides organizations with a structured framework of rules and
procedures for managing sensitive data and guaranteeing its safety and security.

This standard uses a risk-based approach, where organisations continuously monitor security risks
through set policies and procedures, meaning that ISO 27001 provides a comprehensive and
structured framework for protecting information assets and reducing the likelihood of data
breaches and cyber incidents~\cite{ref35}.

One important security measure stated by this standard is to always set up event logging and
monitoring to record system activities and provide help in threat detection and response. XDR
systems align closely with this requirement, as they usually collect information of events from
multiple sources, giving analysts a centralized logging system.

\begin{figure}[H]
  \centering
   \includegraphics[width=0.7\textwidth]{fig16_iso_xdr.png}
  \caption{Conceptual relationship between ISO/IEC 27001 and XDR systems (generated using AI
  tools)}
  \label{fig:iso_xdr}
\end{figure}

% ──────────────────────────────────────────────────────────────────────────────
\section{Compliance Challenges and Operational Limitations in XDR Systems}
% ──────────────────────────────────────────────────────────────────────────────

As mentioned before, XDR systems are the next line of innovation. Along with the enhancements
from AI and ML, as well as the centralization of data, they will be a huge service requested by
all organizations. However, much like the traditional systems currently used in most
cybersecurity solutions (SIEMs, EDRs, SOAR, etc.), they must fully comply with well-known
frameworks and best practices. This might be a challenge, considering that AI tends to be
inconsistent with decision-making, where the algorithms that give AI the ability of learning and
operating autonomously are not yet fine-tuned to fit the steep requirements that the cybersecurity
field has.

% ──────────────────────────────────────────────
\subsection{Transparency and Auditability Issues}
% ──────────────────────────────────────────────

As mentioned before, AI-enhanced XDR systems work like a black box, where the analysts have
difficulties interpreting why the system responds a specific way to a specific incident or
occurrence. The decision process of these tools is impossible to interpret~\cite{ref36}.

\begin{figure}[htbp]
  \centering
   \includegraphics[width=0.75\textwidth]{fig17_blackbox_xai.png}
  \caption{Black Box Concept vs.\ Explainable AI~\cite{ref40}}
  \label{fig:blackbox_xai}
\end{figure}

These issues for auditability and transparency cause big problems, as cybersecurity analysts need
to know and understand why a particular event is deemed malicious, and why a specific incident
response action is taken against it, in order to fully comply with regulations~\cite{ref37}. This
is also an issue for organisations that might be facing an APT (Advanced Persistent Threat),
where the threat is not easily found, making it hard for the analysts to follow up with
understanding why the AI and ML-enhanced XDR system decided to mark a series of activities as
suspicious.

\begin{figure}[htbp]
  \centering
   \includegraphics[width=0.7\textwidth]{fig18_apt_datasets.png}
  \caption{Datasets used for APT detection~\cite{ref37}}
  \label{fig:apt_datasets}
\end{figure}

% ──────────────────────────────────────────────
\subsection{Operational Complexity and Integration}
% ──────────────────────────────────────────────

However, even with the use of AI and strong ML algorithms, analysts can never be one hundred
percent sure that everything that the system deems appropriate is actually appropriate. This begs
the question: with centralized data coming from multiple points of an organization's
infrastructure, how can an administrator handle all this information? They could be rapidly
overwhelmed by the data and have difficulty correlating the events~\cite{ref37}.

Not only that, but the policies enforced should always be firm and consistent, and the response
plan should be standardized, making it even harder for a human analyst to keep up. These issues
become even more problematic when having to deal with largely scaled security solutions that
integrate multiple tools together.

% ──────────────────────────────────────────────
\subsection{Limitations in Automated Compliance Enforcement}
% ──────────────────────────────────────────────

Although the above-mentioned frameworks, such as MITRE ATT\&CK and NIST CSF provide structured
and clear guidance on how to operate safely, the XDR system does not automatically comply with
them.

The MITRE ATT\&CK provides a series of TTPs for well-known cyber-attacks, but that does not mean
that an XDR system should have it as a predefined rule to function a certain way for both
detection and response, thus making MITRE ATT\&CK a reference, and not a main source of
automated compliance.

The same thing can be said about NIST CSF: even though it gives security providers a clear way
of defending an organization's system, it does not mean that the steps will be followed
rigorously using enforced automation, but only through complex organization and
planning~\cite{ref39}. Therefore, an XDR system, even the innovative ones including LLMs and
powerful ML algorithms, still requires manual intervention in some cases.

% ──────────────────────────────────────────────────────────────────────────────
\section{Summary \& Review of Gaps and Future Research Scopes}
% ──────────────────────────────────────────────────────────────────────────────

This chapter provided a review of every modern technology related to cybersecurity systems,
especially focusing on threat and intrusion detection as well as incident response, which are the
two core functionalities of XDR systems, while also provisioning the capabilities and innovations
of these systems incorporated with Large Language Models (LLMs).

It started with traditional and conventional systems, such as SIEMs and EDRs, focusing on their
way of operation as well as their positive impact in providing security, but also covered their
current limitations and their disability to evolve along the threats that are being developed
every day.

The next part covered one of the most important concepts of what the cybersecurity field is
trying to achieve: systems that use LLMs, incorporated with Agentic AI and powerful machine
learning algorithms, to create a centralized dashboard for analysts to easily maintain a whole
infrastructure.

Lastly, the chapter also covered the relationship between these XDR systems and cybersecurity
frameworks, highlighting the importance of the protocol and guidelines that these systems should
follow on every scenario, as well as learn and understand how attackers use tactics and procedures
to infiltrate the system and how to follow a specific set of actions in order to neutralize
large-scale threats.

% ══════════════════════════════════════════════════════════════════════════════
\chapter{System Design and Architecture}
% ══════════════════════════════════════════════════════════════════════════════

This chapter contains insights about the design, development and simulation of the proposed
AI-driven XDR security system. The main points include the overall structure of the solution,
how the components work with each other, and describe the tools and technologies used for
creating it.

The solution's aim is to simulate the behaviour of modern XDR (Extended Detection \& Response)
systems, by collecting telemetry from more than one source, analysing generated logs from these
sources, then identifying any weird behaviour or pattern, to finally automate a response plan.
The way this is going to be achieved is by using a centralized system with powerful Machine
Learning algorithms.

% ──────────────────────────────────────────────────────────────────────────────
\section{System Overview}
% ──────────────────────────────────────────────────────────────────────────────

The system's design can easily be described as an AI-driven XDR security system, that can not
only do what an XDR system alone always does---such as collecting data from multiple sources and
centralizing the display of these data---but also detect suspicious behaviour over different data
layers, correlate this behaviour from one data layer to another, and then respond to this issue
autonomously with the appropriate response plan.

All the components of this simulated solution work with each other, all handling different parts
of the process (detecting or responding), with a main goal of replicating a real XDR system used
in real-world cybersecurity environments. One of the ready-made components that will have to be
included is an open-source telemetry system called Wazuh, alongside a trained AI model, capable
of detecting abnormal patterns and concluding which response action is best appropriate for the
threat.

\begin{figure}[htbp]
  \centering
  \includegraphics[width=0.85\textwidth,height=0.7\textheight,keepaspectratio]{fig19_xdr_system_design.png}
  \caption{XDR Simulated System Design}
  \label{fig:xdr_system_design}
\end{figure}

The above picture gives a visual representation of the layered design, starting from the
endpoints that are being monitored, continuing with Wazuh agents grabbing these logs and sending
them to the detection engine, where powerful Machine Learning algorithms try to detect
abnormalities and identify any possible threats. If there is any suspicious behaviour detected,
the decision engine powered by an AI model determines the severity of the event and simulates a
response plan. The response plan includes things such as process termination, endpoint isolation
and IP blocking.

% ──────────────────────────────────────────────────────────────────────────────
\section{Endpoint Monitoring and Log Collection Layer}
% ──────────────────────────────────────────────────────────────────────────────

The first step of the proposed XDR system is endpoint monitoring and log collection. This layer
is used for simulating the generation of telemetry, collecting it and then sending it to a
centralized processing layer.

The main components are the Wazuh agents, which are installed on the monitored endpoints (in this
case Virtual Machines). They collect data and logs from events and processes and then send them
over to the Wazuh manager for further analysis.

The collected telemetry data can include normal process logging, but also:

\begin{itemize}
  \item Authentication events (not in normal work hours)
  \item Failed login attempts (suspicious if the user has tried multiple times)
  \item Suspicious PowerShell activity (such as the command \texttt{whoami}, to understand user
  access and privileges)
  \item File modification activity (suspicious if a process tries to modify files that are core
  to the operating system, or that include batch processes)
\end{itemize}

\begin{figure}[htbp]
  \centering
   \includegraphics[width=0.5\textwidth]{fig20_whoami.png}
  \caption{\texttt{whoami} Command in PowerShell}
  \label{fig:whoami}
\end{figure}

\begin{figure}[htbp]
  \centering
   \includegraphics[width=0.85\textwidth]{fig21_wazuh_dashboard.png}
  \caption{Wazuh Central Server (Dashboard), accessed on May 23, 2026}
  \label{fig:wazuh_dashboard}
\end{figure}

These processes collected by Wazuh agents, after being sent to the main server, are normalized
and categorized based on already defined rules of category and severity. This gives the AI engine
more capability to understand whether any of these logs represent an abnormality, or if there is
a connection between two or multiple logs that may indicate a possible threat.

% ──────────────────────────────────────────────────────────────────────────────
\section{Central Data Processing and Telemetry Management Layer}
% ──────────────────────────────────────────────────────────────────────────────

This section covers the second part of the solution, specifically talking about how the data taken
from numerous sources is connected to the AI engine which will actually analyse the telemetry.
Anything collected by the Wazuh agents will be sent here, so that the machine learning trained
models can analyse and search for patterns or abnormalities through a wide number of generated
logs.

All the generated telemetry is collected and processed by Wazuh Manager (centralized server
storing all the data), to then be sent through the central processing layer. The most important
thing to do at this stage is to create a standard format for all the logs, which will help the
models better understand the telemetry. This will be done by normalization and use of data
sanitization methods.

The desired end-product will be a prepared dataset that will be further analysed, but there will
also be another crucial step: feature extraction \& labelling. Machine Learning algorithms are
unable to just train themselves on raw telemetry or log data, meaning that specific attributes
will have to be labelled as important variables in the detection and decision-making parts. Table
\ref{tab:feature_labels} shows a raw example of how labelling is made.

\begin{table}[htbp]
  \centering
  \caption{Feature Labelling Example}
  \label{tab:feature_labels}
  \begin{tabular}{>{\bfseries}p{4.5cm} p{9cm}}
    \toprule
    Feature & Description \\
    \midrule
    Failed Login Count    & Suspicious number of login attempts \\
    Event Frequency       & Similar events with same interval apart \\
    Source IP Address     & Activity origin \\
    Process Name          & Name of executed process \\
    Endpoint Severity     & Determined by the Wazuh Manager \\
    Endpoint Identifier   & Finding the source of telemetry \\
    \bottomrule
  \end{tabular}
\end{table}

The design choice of having a different layer for data processing and detection logic followed by
response is applied to current XDR-based cybersecurity solutions, due to the scalability potential
it has for future modifications or additions to the system or to the machine learning algorithms,
also making it easier to just add new endpoints.

\begin{figure}[htbp]
  \centering
   \includegraphics[width=0.8\textwidth]{fig22_central_data_processing.png}
  \caption{Central Data Processing Workflow}
  \label{fig:central_data_processing}
\end{figure}

% ──────────────────────────────────────────────────────────────────────────────
\section{Artificial Intelligence Detection Layer}
% ──────────────────────────────────────────────────────────────────────────────

This is the core part of the XDR solution, offering detailed analysis of all the normalized data
and telemetry received from the Wazuh Manager Server and any real-time logs, helping identify any
discrepancies or abnormalities in a log or a series of logs. Unlike conventional systems, this
will rely not only on basic predefined rules, but also on behavioural analysis where the most
innovative part is the detection of threats that don't use a known pattern of attacking.

The evaluation of the normalized telemetry data is done through powerful Machine Learning
Algorithms, focused on extracting features such as:

\begin{itemize}
  \item Frequency of a process being executed
  \item Network behaviour
  \item Process execution and resource patterns
  \item etc.
\end{itemize}

These features are selected since they determine if a particular event or a thread of executed
processes is of risk or not. The detection phase begins with taking all the extracted features,
sending them to the Machine Learning model, which will classify all the data using an anomaly
score, where the higher the score, the higher the chance of this event or log being an indicator
of a cyber threat.

The design with two different layers---where one classifies potential risks and the other
correlates them---makes it easier to scale the models or the layers themselves.

\begin{figure}[H]
  \centering
   \includegraphics[width=0.8\textwidth,height=0.7\textheight,keepaspectratio]{fig23_ai_detection_layer.png}
  \caption{AI Detection Layer Workflow}
  \label{fig:ai_detection_layer}
\end{figure}

% ──────────────────────────────────────────────────────────────────────────────
\section{Event Correlation and Threat Analysis}
% ──────────────────────────────────────────────────────────────────────────────

The correlation layer is right after the detection layer decides on what to classify as potential
threats, by filtering the suspicious telemetry data out and sending it to this layer. The
correlation logic will evaluate the risks using metrics such as:

\begin{itemize}
  \item \textbf{Threat Scoring:} The higher the score, the higher the chance of a log or event
  being related to a cyber threat.
  \item \textbf{Attack Pattern Matching:} Known sequences of scaled attacks and TTPs (using
  MITRE ATT\&CK's framework) are used to determine whether a sequence of logs matches a
  well-known attack.
  \item \textbf{Confidence Evaluation:} Verifying whether one of the filtered logs can actually
  pose a threat or is just a false positive.
\end{itemize}

An important objective of threat correlation is to identify if there is a connection between
different logs and tie it to a largely scaled attack or intrusion. However, it should also be
able to detect whether the filtered-out logs by the detection layer can be coincidences or false
positives. Through continuous learning and training, better data normalization and feature
extraction, as well as model improvement, it can help create a stronger and more accurate layer.

\begin{figure}[htbp]
  \centering
   \includegraphics[width=0.8\textwidth]{fig24_correlation_layer.png}
  \caption{Correlation Layer Architecture}
  \label{fig:correlation_layer}
\end{figure}

% ──────────────────────────────────────────────────────────────────────────────
\section{Decision and Threat Mitigation Layer}
% ──────────────────────────────────────────────────────────────────────────────

The Decision layer is responsible for validating the results that the AI model has come up with
after suspicious log detection and event correlation. After evaluating the results, it determines
the appropriate response action, acting as the rational thinking and decision-making engine, which
is required in real-world XDR solutions. This feature must be if not fully, mostly automated,
requiring little to no human input.

The automated process is executed based on some performative features, such as:

\begin{itemize}
  \item \textbf{Severity} (based on the score assigned from detection)
  \item \textbf{Confidence} (based on how confident the detection model is that the filtered
  suspicious log is actually an anomaly)
  \item \textbf{Context of detected incident} (The initialization of the attack and if there is
  usage of TTPs)
\end{itemize}

Based on previously generated data such as anomaly scores, correlated events, threat
classification and affected assets, the model continues with assessing the level of risk. Using
this, the model then assigns priority and launches a specific response plan, with examples of
responses below:

\begin{itemize}
  \item Security alert generation
  \item Increased watch and monitor
  \item Blocking of suspicious external sources and IP addresses
  \item Termination of suspicious processes
  \item Isolation of affected assets \& endpoints
\end{itemize}

However, this response action is not directly sent to the affected endpoints, but actually goes
through a testing layer called the Response Simulation Layer, where the chosen response action
is simulated to prove the effectiveness of the choice made by the model.

\begin{figure}[H]
  \centering
   \includegraphics[width=0.8\textwidth, height=0.7\textheight,keepaspectratio]{fig25_decision_layer.png}
  \caption{Decision and Threat Mitigation Layer}
  \label{fig:decision_layer}
\end{figure}

% ──────────────────────────────────────────────────────────────────────────────
\section{Response Simulation Layer}
% ──────────────────────────────────────────────────────────────────────────────

The Response Simulation Layer is the last component that makes this architecture complete. The
response chosen by the AI \& ML model needs to be executed in a controlled testing environment to
ensure that the response is actually effective against the threat. This will replicate the actual
real-world response actions that are executed in any security solution.

The response actions are comprised of actions meant to limit an endpoint's access to the local
network, or even total isolation, thus preventing any malware from spreading to other endpoints
in the same network, stopping ransomware from encrypting files and so on. It is also important
that every action performed by this simulation is saved and logged.

All the logged data from an incident must be done cleanly, making sure that the logging process
is standardized and every generated piece of information can easily be read by a human, or turned
into further training data for future models. Important features include:

\begin{itemize}
  \item Detected Threat
  \item Severity Level
  \item Selected Response Action
  \item Timestamp
  \item Affected Endpoint/Device
  \item Response Status
\end{itemize}

\begin{figure}[htbp]
  \centering
   \includegraphics[width=0.8\textwidth]{fig26_response_execution.png}
  \caption{Response Plan Execution Workflow}
  \label{fig:response_execution}
\end{figure}

% ──────────────────────────────────────────────────────────────────────────────
\section{Dashboard and Monitoring Interface}
% ──────────────────────────────────────────────────────────────────────────────

The dashboard and monitoring interface will serve as the view that the user (analysts) will
interact with the proposed solution. Its stack will comprise React \& TypeScript, simulated by
real-world XDR solutions, that will help show current logs being sent from the different endpoints
in a user-friendly structure.

Besides the data, this dashboard will also preview any incoming security events or active alerts
that are generated at the moment when data is going through the detection and correlation layer,
as well as the results of AI analysis done on the backend. It will also contain a tab where there
will be a chronological order of the issues, giving analysts a detailed view of how the model
processed the information and gave back its response.

\begin{figure}[htbp]
  \centering
   \includegraphics[width=0.9\textwidth]{fig27_dashboard.png}
  \caption{Dashboard Pages}
  \label{fig:dashboard}
\end{figure}

% ──────────────────────────────────────────────────────────────────────────────
\section{System Modelling: UML Diagrams}
% ──────────────────────────────────────────────────────────────────────────────

The design and implementation of the proposed solution require not only functional explanation
and overview, but also a clear correlation between the interactions of users (analysts) and the
system itself. The selected diagrams that are able to fully explain how the system works and how
users interact with it are:

\begin{itemize}
  \item Use Case Diagram
  \item Sequence Diagram
  \item Component Diagram
\end{itemize}

% ──────────────────────────────────────────────
\subsection{Use Case Diagram}
% ──────────────────────────────────────────────

\begin{figure}[htbp]
  \centering
   \includegraphics[width=0.8\textwidth]{fig_usecase.png}
  \caption{Use Case Diagram of AI-driven XDR System}
  \label{fig:use_case}
\end{figure}

The Use Case Diagram is used to illustrate the connection and interaction between system users
and the XDR solution. It helps us understand who are the types of users that can interact with
the system, and what kind of functionality and privilege levels each group of users has. It also
shows how the system actions are connected to each other and can be executed by the users
themselves.

% ──────────────────────────────────────────────
\subsection{Sequence Diagram}
% ──────────────────────────────────────────────

The Sequence Diagram will illustrate not only the interaction between users and components of the
system, but will also give a clear timeline of the execution of each process that the system has,
in a chronological manner. The process will begin from the endpoints that generate logs and event
alerts, which are then collected by Wazuh agents and sent for further inspection. After the data
received is normalized and standardized, it goes through three different layers and displays the
result on the dashboard.

The components that will be part of the diagram are:

\begin{itemize}
  \item Monitored Endpoint
  \item Wazuh Agents
  \item Data Processing Layer
  \item AI Detection Layer
  \item Correlation Layer
  \item Decision-Making Engine
  \item Response Simulation Layer
  \item Dashboard
\end{itemize}

\begin{figure}[htbp]
  \centering
   \includegraphics[width=0.9\textwidth]{fig_sequence.png}
  \caption{Sequence Diagram for XDR solution}
  \label{fig:sequence}
\end{figure}

% ──────────────────────────────────────────────
\subsection{Component Diagram}
% ──────────────────────────────────────────────

The Component Diagram gives an overview of the structural organization of the proposed XDR
solution. Different from the previous diagram that explains the timeline of the interaction
between the user and the components, this diagram provides the internal architecture, including
the data processing module, AI detection and correlation layers, decision engine and the user
dashboard. It also shows the fact that all these components work autonomously, giving the option
to modify or change parts of it without worrying about the whole architecture.

\begin{figure}[H]
  \centering
  \includegraphics[width=0.9\textwidth, height=0.4\textheight,keepaspectratio]{fig_component.png}
  \caption{XDR Solution Component Diagram}
  \label{fig:component}
\end{figure}

% ══════════════════════════════════════════════════════════════════════════════
% References
% ══════════════════════════════════════════════════════════════════════════════
\begin{thebibliography}{40}

\bibitem{ref1}
MITRE ATT\&CK, ``Phishing (T1566),'' \textit{MITRE ATT\&CK\textregistered}.
[Online]. Available: \url{https://attack.mitre.org/techniques/T1566/}
[Accessed: Mar.\ 16, 2026].

\bibitem{ref2}
National Institute of Standards and Technology, \textit{Computer Security Incident Handling
Guide}, NIST Special Publication 800-61r3. Available:
\url{https://doi.org/10.6028/NIST.SP.800-61r3}

\bibitem{ref3}
\url{https://doi.org/10.63282/3050-9262.IJAIDSML-V6I1P111}

\bibitem{ref4}
N.\ N.\ A.\ Sjarif, S.\ Chuprat, M.\ N.\ Mahrin, N.\ A.\ Ahmad, A.\ Ariffin, and
F.\ M.\ Senan, ``Endpoint Detection and Response: Why Use Machine Learning?,''
\textit{IEEE Xplore}, DOI: 10.1109/ICOIN50798.2020.8939836.

\bibitem{ref5}
Microsoft, ``What is XDR?,'' \textit{Microsoft Security}.
[Online]. Available:
\url{https://www.microsoft.com/en-us/security/business/security-101/what-is-xdr}

\bibitem{ref6}
D.\ Pissanidis and K.\ Demertzis, ``Integrating AI/ML in Cybersecurity: An Analysis of Open
XDR Technology and Its Application in Intrusion Detection and System Log Management,''
\textit{Preprints}, Dec.\ 2023, doi: 10.20944/preprints202312.0205.v1.

\bibitem{ref7}
MITRE ATT\&CK, ``MITRE ATT\&CK Framework,'' \textit{MITRE ATT\&CK\textregistered}.
[Online]. Available: \url{https://attack.mitre.org/}

\bibitem{ref8}
N.\ Mohamed, ``Artificial intelligence and machine learning in cybersecurity: a deep dive into
state-of-the-art techniques and future paradigms,'' \textit{Knowledge and Information Systems},
vol.\ 67, pp.\ 6969--7055, 2025. \url{https://doi.org/10.1007/s10115-025-02429-y}

\bibitem{ref9}
S.\ Survaiya, V.\ Uke, and I.\ Vighe, ``AI in cybersecurity: Intrusion detection system,''
\textit{International Advanced Research Journal in Science, Engineering and Technology},
vol.\ 12, no.\ 11, 2025, doi: 10.17148/IARJSET.2025.121106.

\bibitem{ref10}
N.\ Mohamed, ``Current trends in AI and ML for cybersecurity: A state-of-the-art survey,''
\textit{Cogent Engineering}, vol.\ 10, no.\ 2, p.\ 2272358, 2023,
DOI: 10.1080/23311916.2023.2272358.

\bibitem{ref11}
S.\ Survaiya, V.\ Uke, and I.\ Vighe, ``AI in cybersecurity: Intrusion detection system,''
\textit{International Advanced Research Journal in Science, Engineering and Technology},
vol.\ 12, no.\ 11, Nov.\ 2025, doi: 10.17148/IARJSET.2025.121106.

\bibitem{ref12}
D.\ Tomar, K.\ Chhillar, and D.\ Kumhar, ``Comparative study of supervised, unsupervised and
reinforcement learning approaches for malware detection,'' \textit{International Journal of
Innovative Research in Technology (IJIRT)}, vol.\ 12, no.\ 4, pp.\ 2793--2801, Sep.\ 2025.

\bibitem{ref13}
Canadian Institute for Cybersecurity, ``CICIDS2017 Dataset.'' [Online]. Available:
\url{https://www.unb.ca/cic/datasets/ids-2017.html}

\bibitem{ref14}
S.\ T.\ Erukude, V.\ C.\ Marella, and S.\ R.\ Veluru, ``AI-Driven Cybersecurity Threats:
A Survey of Emerging Risks and Defensive Strategies,'' arXiv:2601.03304, Jan.\ 2026.
[Online]. Available: \url{https://arxiv.org/abs/2601.03304}

\bibitem{ref15}
J.\ Pan, ``AI based Log Analyser: A Practical Approach,'' arXiv, 2022.

\bibitem{ref16}
NSA \& CISA, ``Implementing SIEM and SOAR Platforms Guidance,'' 2025.

\bibitem{ref17}
A.\ Ehis, ``Optimization of SIEM Infrastructures and Event Correlation Analysis,'' 2023.

\bibitem{ref18}
N.\ Sharma and N.\ Jindal, ``Emerging artificial intelligence applications: metaverse, IoT,
cybersecurity, healthcare -- an overview,'' \textit{Multimedia Tools and Applications},
vol.\ 83, no.\ 19, pp.\ 57317--57345, 2024.

\bibitem{ref19_fig}
C.\ Kyrkou et al., ``Towards artificial-intelligence-based cybersecurity for robustifying
automated driving systems against camera sensor attacks,'' in \textit{Proc.\ 2020 IEEE
Computer Society Annual Symposium on VLSI (ISVLSI)}, 2020.

\bibitem{ref19b}
C.\ Kyrkou et al., ``Towards artificial-intelligence-based cybersecurity for robustifying
automated driving systems against camera sensor attacks,'' in \textit{Proc.\ 2020 IEEE
Computer Society Annual Symposium on VLSI (ISVLSI)}, 2020.

\bibitem{ref20}
S.\ Dua and X.\ Du, ``Data Mining and Machine Learning in Cybersecurity,''
arXiv:1904.03491, 2019. [Online]. Available: \url{https://arxiv.org/abs/1904.03491}

\bibitem{ref21}
N.\ Mohamed, ``Current Trends in AI and ML for Cybersecurity: A State-of-the-Art Survey,''
\textit{Cogent Engineering}, 2023.

\bibitem{ref22}
I.\ Kim et al., ``Toward Robust Security Orchestration and Automated Response,''
\textit{Information}, 2025.

\bibitem{ref23}
M.\ S.\ Arsanjani, ``The Anatomy of Agentic AI,'' \textit{Medium}, 2023. Available:
\url{https://dr-arsanjani.medium.com/the-anatomy-of-agentic-ai-0ae7d243d13c}
[Accessed: May 19, 2026].

\bibitem{ref24}
N.\ Katiyar, ``AI And Cyber-Security: Enhancing Threat Detection And Response With Machine
Learning,'' \textit{Educational Administration: Theory and Practice}, vol.\ 30, no.\ 4,
pp.\ 6273--6282, 2024. doi: 10.53555/kuey.v30i4.2377.

\bibitem{ref25}
G.\ Palma, G.\ Cecchi, M.\ Caronna, and A.\ Rizzo, ``Leveraging Large Language Models for
Scalable and Explainable Cybersecurity Log Analysis,'' \textit{Journal of Cybersecurity and
Privacy}, vol.\ 5, p.\ 55, 2025. \url{https://doi.org/10.3390/jcp5030055}

\bibitem{ref26}
Skedler, ``What is Explainable AI (XAI)?'' 2023. [Online]. Available:
\url{https://www.skedler.com/blog/what-is-explainable-ai-xai/}
[Accessed: May 19, 2026].

\bibitem{ref27}
A.\ Vassilev, A.\ Oprea, A.\ Fordyce, and H.\ Anderson, \textit{Adversarial Machine Learning:
A Taxonomy and Terminology of Attacks and Mitigations}, NIST AI 100-2e2023, 2024.
\url{https://doi.org/10.6028/NIST.AI.100-2e2023}

\bibitem{ref28}
Simplilearn, ``CIA Triad in Cyber Security,'' 2026. [Online]. Available:
\url{https://www.simplilearn.com/cia-triad-in-cyber-security-article}
[Accessed: May 19, 2026].

\bibitem{ref29}
G.\ Sarraf and V.\ Pal, ``Autonomous Threat Detection and Response in Cloud Security:
A Comprehensive Survey of AI-Driven Strategies,'' 2025.

\bibitem{ref30}
V.\ Shewale, ``Beyond EDR: Exploring the rise of XDR for unified threat detection and
response,'' \textit{World Journal of Advanced Engineering Technology and Sciences},
vol.\ 15, no.\ 02, pp.\ 380--386, 2025.
doi: \url{https://doi.org/10.30574/wjaets.2025.15.2.0551}

\bibitem{ref31}
Microsoft, ``What is the MITRE ATT\&CK framework?,'' \textit{Microsoft Security}.
\url{https://www.microsoft.com/en-us/security/business/security-101/what-is-mitre-attack-framework}
[Accessed: May 20, 2026].

\bibitem{ref32}
IBM, ``What is the MITRE ATT\&CK framework?,'' \textit{IBM Think}.
\url{https://www.ibm.com/think/topics/mitre-attack} [Accessed: May 20, 2026].

\bibitem{ref33}
MITRE Corporation, ``ATT\&CK,'' \textit{MITRE ATT\&CK}. \url{https://attack.mitre.org/}

\bibitem{ref34}
National Institute of Standards and Technology (NIST), ``Identify, Protect, Detect, Respond
and Recover: The NIST Cybersecurity Framework,'' \textit{NIST Taking Measure Blog}.
\url{https://www.nist.gov/blogs/taking-measure/identify-protect-detect-respond-and-recover-nist-cybersecurity-framework}
[Accessed: May 20, 2026].

\bibitem{ref35}
IBM, ``ISO 27001 compliance -- IBM Cloud,'' \textit{IBM Products}.
\url{https://www.ibm.com/products/cloud/compliance/iso-27001}
[Accessed: May 20, 2026].

\bibitem{ref36}
X.\ Li et al., ``Explainable dimensionality reduction (XDR) to unbox AI `black box' models,''
\textit{Humanities and Social Sciences Communications}, vol.\ 10, p.\ 35, 2023.
\url{https://doi.org/10.1057/s41599-023-01505-4}

\bibitem{ref37}
N.\ H.\ A.\ Mutalib et al., ``Explainable deep learning approach for advanced persistent
threats (APTs) detection in cybersecurity: a review,'' \textit{Artificial Intelligence Review},
vol.\ 57, p.\ 297, 2024. \url{https://doi.org/10.1007/s10462-024-10890-4}

\bibitem{ref38}
A.\ S.\ George, A.\ George, Dr.\ Baskar, and D.\ Pandey, ``XDR: The Evolution of Endpoint
Security Solutions,'' \textit{International Journal of Advanced Research in Science
Communication and Technology}, vol.\ 8, pp.\ 493--501, 2021.
doi: 10.5281/zenodo.7028219.

\bibitem{ref39}
J.\ Jana, ``The AI Black Box: What We're Still Getting Wrong about Trusting Machine Learning
Models,'' \textit{Hyperight}, Mar.\ 3, 2025. [Online]. Available:
\url{https://hyperight.com/ai-black-box-what-were-still-getting-wrong-about-trusting-machine-learning-models/}
[Accessed: May 11, 2026].

\bibitem{ref40}
A.\ Hui, S.\ Ahn, C.\ Lye, and J.\ Deng, ``Ethical Challenges of Artificial Intelligence in Health Care: A Narrative Review,'' \textit{Ethics in Biology, Engineering and Medicine}, vol.\ 12, 2022. doi: 10.1615/EthicsBiologyEngMed.2022041580.

\end{thebibliography}

\end{document}